% Part: sets-functions-relations
% Chapter: functions
% Section: partial-functions

\documentclass[../../../include/open-logic-section]{subfiles}


\begin{document}

\olfileid{sfr}{fun}{par}

\olsection{अंशतः फलने}

\begin{explain}
कधी कधी फलनाच्या व्याख्येत थोडी सवलत देणे उपयोगी ठरते: सर्व शक्य
आदानांसाठी फलनाचे प्रदान परिभाषित असलेच पाहिजे ही अट काढून टाकता
येते. अशा संगतींना \emph{अंशतः फलने} म्हणतात.
\end{explain}

\begin{defn}
\emph{अंशतः फलन} $f \colon A \pto B$ ही अशी संगती असते की ती~$A$
मधील प्रत्येक !!{element} शी~$B$ मधील जास्तीत जास्त एक !!{element}
जोडते. $f$ ने $x \in A$ शी~$B$ मधील एखादा घटक जोडला असेल, तर
$f(x)$ \emph{परिभाषित} आहे असे म्हणतो; अन्यथा ते \emph{अपरिभाषित}
असते. $f(x)$ परिभाषित असल्यास $f(x) \fdefined$ असे लिहितो;
अन्यथा $f(x) \fundefined$ असे लिहितो. अंशतः फलन~$f$ चा \emph{प्रांत}
म्हणजे ते परिभाषित असलेल्या~$A$ मधील घटकांचा उपसंच:
$\dom{f} = \Setabs{x \in A}{f(x) \fdefined}$.
\end{defn}

\begin{ex}
प्रत्येक फलन $f\colon A \to B$ हे अंशतः फलनही असते. ~$A$ वरील
सर्व ठिकाणी परिभाषित असलेल्या अंशतः फलनांना---म्हणजे ज्यांना आपण
आतापर्यंत फक्त फलन म्हणत आलो, त्यांना---\emph{सर्वत्र परिभाषित}
फलने असेही म्हणतात.
\end{ex}

\begin{ex}
$f(x) = 1/x$ ने दिलेले अंशतः फलन $f \colon \Real \pto \Real$
हे $x = 0$ साठी अपरिभाषित असते आणि इतर सर्व ठिकाणी परिभाषित असते.
\end{ex}

\begin{prob}
$f\colon A \pto B$ दिले असता, अंशतः फलन $g\colon B \pto A$ ची
व्याख्या अशी करा: कोणत्याही $y \in B$ साठी $f(x) = y$ असणारा एकमेव
$x \in A$ असेल, तर $g(y) = x$; अन्यथा $g(y) \fundefined$.
$f$ हे एकास-एक असल्यास सर्व $x \in \dom{f}$ साठी $g(f(x)) = x$
आणि सर्व $y \in \ran{f}$ साठी $f(g(y)) = y$ असते, हे दाखवा.
\end{prob}

\begin{defn}[अंशतः फलनाचा आलेख]
$f\colon A \pto B$ हे अंशतः फलन असू द्या. ~$f$ चा \emph{आलेख}
म्हणजे पुढील व्याख्येने मिळणारा संबंध $R_f \subseteq A \times B$:
\[
R_f = \Setabs{\tuple{x,y}}{f(x) = y}.
\]
\end{defn}

\begin{prop}
$R \subseteq A \times B$ या संबंधाचा असा गुणधर्म आहे असे समजा की
$Rxy$ आणि $Rxy'$ असल्यास $y = y'$ असते. मग $R$ हा पुढीलप्रमाणे
व्याख्या केलेल्या अंशतः फलन $f\colon A \pto B$ चा आलेख असतो:
$Rxy$ असणारा एक $y$ असेल, तर $f(x) = y$; अन्यथा $f(x) \fundefined$.
शिवाय $R$ हा \emph{प्रत्येक आदानाशी किमान एक प्रदान जोडणारा} असेल,
म्हणजे प्रत्येक $x \in A$ साठी $Rxy$ असणारा एक $y \in B$ असेल,
तर $f$ हे सर्वत्र परिभाषित असते.
\end{prop}

\begin{proof}
$Rxy$ असणारा एक $y$ आहे असे समजा. $Rxy'$ असणारा दुसरा
$y' \neq y$ असता, तर $R$ वरील अट मोडली गेली असती. म्हणून
$Rxy$ असणारा एक $y$ असेल, तर तो $y$ एकमेव असतो. त्यामुळे $f$
सुस्पष्टपणे परिभाषित आहे. उघडच $R_f = R$, आणि~$R$ प्रत्येक आदानाशी
किमान एक प्रदान जोडत असेल, तर $f$ हे सर्वत्र परिभाषित असते.
\end{proof}

\end{document}
